
\begin{obeassessment}
\textbf{Kuis Bab 6}
1. Berikan contoh skenario di mana mode CTR lebih menguntungkan dibandingkan mode CBC.
2. Apa yang dimaksud dengan \textit{error propagation} pada mode CBC dan bagaimana mode OFB mengatasinya?
3. Mengapa penggunaan IV (\textit{Initialization Vector}) sangat penting pada mode CBC?
4. Jelaskan mengapa mode ECB dapat dimanipulasi oleh pihak lawan untuk mengubah makna pesan meskipun kuncinya tidak diketahui.
5. Sebuah cipher blok berukuran 4 bit memakai $E_K(P)=(P\oplus K)\ll 1$ dengan $K=\code{1011}$. Hitung cipherteks dari plainteks $\code{0110}$ pada mode ECB.
6. Pada mode OFB, mengapa aliran kunci dapat dihitung lebih dahulu sebelum plainteks diketahui? Sebutkan satu keuntungan praktis dari sifat itu.
7. Jelaskan hubungan antara jumlah putaran sebuah \textit{iterated cipher} dengan tingkat keamanan dan dengan biaya komputasi. Mengapa jumlah putaran tidak dibuat sebesar mungkin?
8. Apa perbedaan utama antara jaringan Feistel dan jaringan SPN dari sudut pandang perancang algoritma?
\end{obeassessment}
\begin{competencychecklist}
\begin{tabularx}{\linewidth}{|X|c|}
\hline
Indikator Kompetensi & Tercapai \\
\hline
Saya dapat menjelaskan perbedaan \textit{block cipher} dan \textit{stream cipher} & [ ] \\
\hline
Saya dapat menjelaskan mengapa $E_K$ harus merupakan permutasi bijektif & [ ] \\
\hline
Saya dapat menangani blok terakhir yang tidak lengkap dengan \textit{padding} & [ ] \\
\hline
Saya dapat membedakan mode operasi ECB, CBC, CFB, OFB, dan CTR & [ ] \\
\hline
Saya dapat menuliskan persamaan enkripsi dan dekripsi tiap mode operasi & [ ] \\
\hline
Saya dapat menganalisis penjalaran galat pada mode CBC, CFB, dan OFB & [ ] \\
\hline
Saya dapat menjelaskan peran IV pada CBC dan peran counter pada CTR & [ ] \\
\hline
Saya dapat menjelaskan prinsip \textit{confusion} dan \textit{diffusion} & [ ] \\
\hline
Saya dapat menghitung substitusi satu masukan pada S-box DES & [ ] \\
\hline
Saya dapat menjelaskan efek \textit{avalanche} pada cipher blok & [ ] \\
\hline
Saya dapat menguraikan mekanisme kerja Jaringan Feistel & [ ] \\
\hline
Saya dapat membuktikan sifat reversibel Jaringan Feistel & [ ] \\
\hline
Saya dapat menjelaskan tujuan \textit{key scheduling} & [ ] \\
\hline
Saya dapat menganalisis pengaruh panjang kunci terhadap keamanan sistem & [ ] \\
\hline
\end{tabularx}
\end{competencychecklist}
